The implementation does not stop at perception. A central design goal is
to make perception actionable in flight-relevant terms, and this is why
the runtime introduces two distinct but complementary control paths: one
centered on Bitcraze Crazyflie for indoor experimentation and one
centered on PX4-compatible communication through MAVLink for a more
standard professional autopilot ecosystem.

The Crazyflie path is valuable because it gives the research platform an
immediately usable indoor drone target that is lightweight, inexpensive,
and safe to iterate on in constrained spaces. In implementation terms,
the runtime exposes high-level flight functions for arming, takeoff,
landing, attitude control, hover-like motion, waypoint movement, and
basic link validation. The significance of this layer is less in the
underlying protocol details and more in the abstraction boundary it
creates: the experimental logic written in MicroPython can command a
real flying platform without descending into firmware-specific
implementation code for every new experiment.

The MAVLink path plays a different role. It positions the same
perception and scripting environment toward a standard autopilot
workflow, enabling connection to ground-control software and to a
PX4-style vehicle architecture. In the implementation, the MAVLink
bridge is built around the official message library and a dedicated
receive task, with queue-based delivery of decoded telemetry into the
runtime. This allows the platform to consume motion, position, and
status telemetry while also publishing onboard detections and system
messages back to the wider vehicle ecosystem. The result is a credible
bridge between a compact onboard perception unit and a more conventional
autopilot stack.

This bridge is not limited to a single control primitive. The runtime
exposes both transmission and reception semantics at the scripting
layer, including reception of autopilot telemetry and chunked
transmission of vision payloads that would not fit cleanly into a single
small embedded message. That detail matters because it shows the design
is meant for sustained exchange between perception and flight software
rather than for one-shot demo messages.

This dual-target approach is important conceptually. Many embedded AI
demonstrators either remain tied to a single experimental robot or aim
immediately at a professional stack without a low-risk iteration
platform. Here, the implementation supports both modes. Crazyflie serves
as a rapid prototyping platform for indoor validation, while MAVLink
provides compatibility with broader drone workflows. From a research
standpoint, this strengthens the portability of the proposed autonomy
layer.

The communication stack extends further through mesh and backend
services. The firmware side includes a compact binary vision message
protocol and bridges for wireless dissemination, while the repository
also adds backend proxy services that relay, aggregate, and expose
telemetry through ground-side interfaces such as MQTT. This is
significant because it shows that the implementation is intended to
operate as part of a distributed autonomy pipeline rather than as an
isolated board demo. Detections are not merely computed; they are
structured, transported, and made available to external systems.

In implementation terms, the control and communication layers provide
the missing operational dimension of embedded perception. A detector
running locally is useful only to a point. A detector that can feed a
Crazyflie experiment, inform a MAVLink autopilot context, publish
compact messages to external services, and remain scriptable from a
high-level runtime becomes much more valuable as a research instrument.
This is the level at which the runtime begins to resemble a true onboard
autonomy substrate rather than a perception benchmark harness.
